\problem{2.29}{Odd bound states of a finite square well}
For $V(x)=-V_0$ on $|x|<a$ and $0$ outside,
bound levels satisfy $-V_0<E<0$. Define
\begin{equation}
l=\frac{\sqrt{2m(E+V_0)}}{\hbar},\quad
\kappa=\frac{\sqrt{-2mE}}{\hbar},\quad
z=la,\quad z_0=\frac{a\sqrt{2mV_0}}{\hbar}.
\end{equation}
An odd eigenfunction is $D\sin(lx)$ inside and
$F\operatorname{sgn}(x)\ee^{-\kappa|x|}$ outside.
At $x=a$,
\begin{equation}
F\ee^{-\kappa a}=D\sin(la),\qquad
-\kappa F\ee^{-\kappa a}=lD\cos(la).
\end{equation}
Dividing yields $l\cot(la)=-\kappa$, or
\begin{equation}
\boxed{-z\cot z=\sqrt{z_0^2-z^2}},\qquad 0<z<z_0.
\end{equation}
This equation can be solved by intersections of
$-z\cot z$ with $\sqrt{z_0^2-z^2}$ on the
intervals where the left side is positive.
For a deep well ($z_0\gg1$), roots approach
$z=n\pi/2$ for even $n$, corresponding to
odd-parity states of the centered infinite well.
As the well weakens the first odd state appears
only when $z_0>\pi/2$; there is no odd bound
state for arbitrarily weak $V_0$.
The even ground state, by contrast, always exists.

\problem{2.30}{Normalizing a finite-well even bound state}
The even solution is
\begin{equation}
\psi(x)=
\begin{cases}
D\cos(lx),&|x|<a,\\
F\ee^{-\kappa|x|},&|x|\ge a.
\end{cases}
\end{equation}
Continuity at $x=a$ gives
$F\ee^{-\kappa a}=D\cos(la)$,
or $F=D\cos(la)\ee^{\kappa a}$.
Hence
\begin{align}
1&=2D^2\int_0^a\cos^2(lx)\dd x+
2F^2\int_a^\infty\ee^{-2\kappa x}\dd x\\
&=D^2\left[a+\frac{\sin(2la)}{2l}
+\frac{\cos^2(la)}{\kappa}\right].
\end{align}
Choose the overall phase so $D>0$:
\begin{equation}
\boxed{D=\left[a+\frac{\sin(2la)}{2l}
+\frac{\cos^2(la)}{\kappa}\right]^{-1/2}},
\quad
\boxed{F=D\cos(la)\ee^{\kappa a}}.
\end{equation}
The even-state eigenvalue relation $\kappa=l\tan(la)$
from the textbook must be solved before evaluating
these constants numerically.

\problem{2.31}{The delta well as the zero-width square-well limit}
Let the square well of half-width $a$ have depth
$V_0=\alpha/(2a)$, so its integrated strength
is $-2aV_0=-\alpha$. As $a\to0$,
\begin{equation}
z_0^2=\frac{2mV_0a^2}{\hbar^2}
=\frac{m\alpha a}{\hbar^2}\longrightarrow0.
\end{equation}
For the even bound level $\kappa=l\tan(la)$ and
$la=z\to0$, so
\begin{equation}
\kappa\sim l^2a
=\frac{2m(E+V_0)a}{\hbar^2}
\longrightarrow\frac{m\alpha}{\hbar^2}.
\end{equation}
Therefore
\begin{equation}
\boxed{E\longrightarrow
-\frac{\hbar^2\kappa^2}{2m}
=-\frac{m\alpha^2}{2\hbar^2}},
\end{equation}
the single attractive-delta bound energy.
For positive incident energy,
the square-well transmission formula is
\begin{equation}
T^{-1}=1+\frac{V_0^2}{4E(E+V_0)}
\sin^2\left(\frac{2a\sqrt{2m(E+V_0)}}{\hbar}\right).
\end{equation}
Using $\sin z\sim z$ and $2aV_0=\alpha$,
\begin{equation}
T^{-1}\to1+\frac{m\alpha^2}{2\hbar^2 E}
=1+\left(\frac{m\alpha}{\hbar^2k}\right)^2,
\quad k=\sqrt{2mE}/\hbar.
\end{equation}
The apparent infinite depth is therefore
a weak-well limit in the finite-square-well
dimensionless parameter $z_0$.

\problem{2.32}{Finite square-well scattering}
For $E>0$ define
$k=\sqrt{2mE}/\hbar$ outside and
$l=\sqrt{2m(E+V_0)}/\hbar$ inside
the well $-a<x<a$.
Write the left and right solutions
\begin{equation}
\psi_L=A\ee^{\ii kx}+B\ee^{-\ii kx},\qquad
\psi_R=F\ee^{\ii kx}
\end{equation}
and the interior as
$\psi_I=C\sin(lx)+D\cos(lx)$.
The matching equations for both $\psi$ and
$\psi'$ at $\pm a$ give
\begin{equation}
\frac{A}{F}
=\ee^{2\ii ka}
\left[\cos(2la)-\frac{\ii}{2}
\left(\frac{k}{l}+\frac{l}{k}\right)\sin(2la)\right].
\end{equation}
Taking the squared modulus and using
$k^2-l^2=-2mV_0/\hbar^2$,
\begin{align}
\frac{|A|^2}{|F|^2}
&=\cos^2(2la)+\frac14
\left(\frac{k}{l}+\frac{l}{k}\right)^2\sin^2(2la)\\
&=1+\frac{(l^2-k^2)^2}{4k^2l^2}\sin^2(2la).
\end{align}
Since exterior velocities agree, $T=|F/A|^2$:
\begin{equation}
\boxed{T=
\left[1+\frac{V_0^2}{4E(E+V_0)}
\sin^2\left(\frac{2a\sqrt{2m(E+V_0)}}{\hbar}\right)
\right]^{-1}}.
\end{equation}
Resonances occur when $2la=n\pi$, making $T=1$.

\problem{2.33}{Transmission across a rectangular barrier}
For a barrier $V_0>0$ on $|x|<a$,
its full width is $L=2a$. Repeating the matching
of Problem 2.32 with interior wave number
$q=\sqrt{2m(E-V_0)}/\hbar$ gives for $E>V_0$
\begin{equation}
\boxed{T^{-1}=1+
\frac{V_0^2}{4E(E-V_0)}
\sin^2\left(\frac{L\sqrt{2m(E-V_0)}}{\hbar}\right)}.
\end{equation}
For $E<V_0$ put
$q=\ii\kappa$, $\kappa=\sqrt{2m(V_0-E)}/\hbar$,
and use $\sin(\ii y)=\ii\sinh y$. This yields
\begin{equation}
\boxed{T^{-1}=1+
\frac{V_0^2}{4E(V_0-E)}\sinh^2(\kappa L)}.
\end{equation}
The $E=V_0$ case is the continuous limit:
$\sin(qL)/q\to L$, leading to
\begin{equation}
\boxed{T(E=V_0)=\left[
1+\frac{mV_0 L^2}{2\hbar^2}\right]^{-1}}.
\end{equation}
For $E<V_0$, $T$ is exponentially suppressed at
large $\kappa L$, but is not zero: quantum tunneling.

\problem{2.34}{Reflection and transmission at an upward step}
Let $V=0$ for $x<0$ and $V=V_0$ for $x>0$.
The incident wave has $k=\sqrt{2mE}/\hbar$.
For $0<E<V_0$ the right-hand solution is
$F\ee^{-\kappa x}$ with
$\kappa=\sqrt{2m(V_0-E)}/\hbar$. Continuity of
$\psi$ and $\psi'$ at zero yields
\begin{equation}
\frac BA=\frac{k-\ii\kappa}{k+\ii\kappa},
\qquad \boxed{R=1,\quad T=0}.
\end{equation}
An evanescent tail exists for $x>0$ but carries no
transmitted flux. For $E>V_0$ let
$q=\sqrt{2m(E-V_0)}/\hbar$ and
$\psi_R=F\ee^{\ii qx}$.
Matching $A+B=F$, $k(A-B)=qF$ yields
\begin{equation}
\boxed{\frac BA=\frac{k-q}{k+q},
\quad \frac FA=\frac{2k}{k+q}}.
\end{equation}
The current of a wave $C\ee^{\ii kx}$ is
$\hbar k|C|^2/m$, so
\begin{equation}
\boxed{R=\frac{(k-q)^2}{(k+q)^2}},\qquad
\boxed{T=\frac qk\frac{|F|^2}{|A|^2}
=\frac{4kq}{(k+q)^2}},
\end{equation}
and $R+T=1$. The factor $q/k$ is essential
because the particle speed changes across the step.

\problem{2.35}{Quantum reflection from a downward step}
For $V=0$ on $x<0$ and $V=-V_0$ on $x>0$,
the transmitted wave number is
$q=\sqrt{2m(E+V_0)}/\hbar$.
Matching as before gives
\begin{equation}
R=\left(\frac{k-q}{k+q}\right)^2
=\left(\frac{\sqrt E-\sqrt{E+V_0}}
{\sqrt E+\sqrt{E+V_0}}\right)^2 .
\end{equation}
At $E=V_0/3$, the ratio $q/k=\sqrt{1+3}=2$,
so $\boxed{R=1/9}$ and $\boxed{T=8/9}$.
This also applies to the simplified example of
a $4$ MeV neutron approaching a $12$ MeV nuclear
potential drop: the probability of
\emph{entering} the interior is $8/9$ in
this one-dimensional step model.
\qcheck{Transmission through the nuclear surface
is not identical to irreversible nuclear absorption;
that additional process requires a more complete
interaction model. A falling car at a cliff also does
not experience a sudden change of gravitational
potential energy at the edge; the actual potential
depends on height and changes during the fall.}
